% GENERATED by tools/build.py -- do not edit

\section{Borcea--Br\"and\'en AIM problems}

\subsection{Problems 35 and 38: mixed-determinant representation and Permanent-On-Top}
\casecredits{Posed by Borcea and Br\"and\'en \citeyearpar{AIMStableProblems}; counterexample found by TODO (TODO); formalized by S.~Sra; audited by S.~Sra.}
A polynomial $f\in\R[z_1,\ldots,z_n]$ is \emph{real stable} if $f(z_1,\ldots,z_n)\ne0$ whenever every $z_i$ has strictly positive imaginary part.  A symmetric polynomial that is homogeneous of degree $d$ expands uniquely in the basis of Schur polynomials as $p=\sum_{\lambda\vdash d}a_\lambda s_\lambda$, and $p$ is called \emph{Schur-positive} when every $a_\lambda\ge0$.  For a partition $\lambda\vdash d$ we write $f^\lambda$ for the number of standard Young tableaux of shape $\lambda$, $\lambda'$ for the conjugate partition, and, for a $d\times d$ matrix $A=(a_{ij})$,
\[
\Imm_\lambda(A)=\sum_{\sigma\in S_d}\chi^\lambda(\sigma)\prod_{i=1}^da_{i\sigma(i)}
\]
for the immanant attached to the irreducible character $\chi^\lambda$ of $S_d$; the two extreme cases are $\Imm_{(1^d)}=\det$ and $\Imm_{(d)}=\per$.  Finally, for $n\times n$ matrices $A_1,\ldots,A_m$ the mixed-determinant, or Johnson, polynomial is
\[
\eta(A_1,\ldots,A_m)=\sum_{(S_1,\ldots,S_m)}\det(A_1[S_1])\cdots\det(A_m[S_m]),
\]
the sum running over all ordered partitions of $\{1,\ldots,n\}$ into pairwise disjoint sets $S_1,\ldots,S_m$ with union $\{1,\ldots,n\}$, and $A[S]$ denoting the principal submatrix of $A$ on the rows and columns in $S$.  When $A\succeq0$ the polynomial $\eta(z_1A,\ldots,z_nA)$ is symmetric, homogeneous, and real stable, and it satisfies $\eta(z_1A,\ldots,z_nA)=\sum_{\lambda\vdash n}\Imm_{\lambda'}(A)s_\lambda(z_1,\ldots,z_n)$ \cite{AIMStableProblems}.  The two problems below ask how far those properties go the other way.
\begin{problem}[As posed: Borcea--Br\"and\'en, AIM problem list \citeyearpar{AIMStableProblems}, Problem 35]
Characterize all symmetric, real stable and homogeneous polynomial of degree $d$ in $n$ variables.  Is it in fact so that if $p$ is such a polynomial, then there is a $d\times d$ positive semidefinite $d\times d$ matrix $A$ such that
\[p(z_1,\ldots,z_n)=\eta(z_1A,\ldots,z_nA)?\]
This is not obvious even for $n=2$.
\end{problem}
\begin{problem}[As posed: Borcea--Br\"and\'en, AIM problem list \citeyearpar{AIMStableProblems}, Problem 38]
Lieb's Permanent-On-Top (POT) Conjecture for the symmetric group on $d$ elements [J,L] asserts that if $A$ is a $d\times d$ positive semi-definite matrix then $\Imm_\lambda(A)\le f^\lambda\per(A)$ for all $\lambda\vdash d$.  An affirmative answer to the following problem would in particular imply the validity of the POT Conjecture.

\emph{Problem 38.}  Let $p$ be as in Problem 36.  Is $a_\lambda\le f^\lambda a_{(1^d)}$ for all $\lambda\vdash d$?
\end{problem}
Borcea--Br\"and\'en posed the relevant representation and endpoint problems in the AIM problem list \citeyearpar{AIMStableProblems}.  Let $S=x_1+\cdots+x_5$ and
\begin{equation}\label{eq:stable-p}
p(x_1,\ldots,x_5)=\prod_{i=1}^5\left(x_i+5\sum_{j\ne i}x_j\right)=\prod_{i=1}^5(5S-4x_i).
\end{equation}
Every factor has strictly positive coefficients, hence has positive imaginary part whenever all variables do; therefore $p$ is real stable.  It is symmetric, homogeneous of degree five in five variables, and has strictly positive monomial coefficients.

\begin{theorem}[\statusfalse: Stable Permanent-On-Top, AIM Problem 38]\label{thm:pot}
The polynomial \eqref{eq:stable-p} satisfies all hypotheses of Borcea--Br\"and\'en Problem 36 but violates
\[
a_\lambda\le f^\lambda a_{(1^5)}.
\]
\end{theorem}
\begin{proof}
Its exact Schur expansion is
\begin{align*}
p={}&625s_{(5)}+4500s_{(4,1)}+7125s_{(3,2)}+8150s_{(3,1,1)}\\
&+7525s_{(2,2,1)}+6580s_{(2,1,1,1)}+1281s_{(1^5)}.
\end{align*}
For $\lambda=(3,2)$, the hook-length formula gives $f^{(3,2)}=5$, whereas
\[
7125>5\cdot1281=6405.
\]
The same upper endpoint fails for $(3,1,1)$, $(2,2,1)$, and $(2,1,1,1)$.  The source package independently reconstructs the monomial expansion and inverts the Kostka matrix over $\mathbb Z$.
\end{proof}

The polynomial also has a diagonal positive-definite mixed-determinant representation.  For $j=1,\ldots,5$, let
\[
D_j=\diag(5,\ldots,5,\underset{j}{1},5,\ldots,5).
\]
Then the multivariate Johnson polynomial satisfies
\[
\eta(x_1D_1,\ldots,x_5D_5)=p(x_1,\ldots,x_5).
\]

\begin{theorem}[\statusfalse: common-matrix representation, AIM Problem 35]\label{thm:common-A}
There is no $5\times5$ Hermitian positive semidefinite matrix $A$ such that
\[
p(x_1,\ldots,x_5)=\eta(x_1A,\ldots,x_5A).
\]
\end{theorem}
\begin{proof}
Such a representation would identify the Schur coefficients, up to one common positive normalization, with irreducible immanants of $A$ indexed by conjugate partitions.  Immanant Permanent-On-Top is known in size five, indeed in every size at most thirteen \cite{Wanless2022}.  It would therefore force
$a_\lambda\le f^\lambda a_{(1^5)}$ for every $\lambda\vdash5$, contradicting Theorem~\ref{thm:pot}.
\end{proof}

\begin{remark}[Scope]
The example is Schur-positive, and its lower endpoint inequalities hold, so it refutes neither Problem 36 nor Problem 37; those require the positive-definite pencil of Theorem~\ref{thm:aim36} below.  The full independent audit is included as \path{counterexamples/stable-schur/artifacts/}.
\end{remark}

\subsection{Problem 36: stable Schur positivity}
\casecredits{Posed by Borcea and Br\"and\'en \citeyearpar{AIMStableProblems}; counterexample found by TODO (TODO); formalized by S.~Sra; audited by S.~Sra.}
We keep the notation of the previous subsection: real stability means nonvanishing when every variable lies in the open upper half-plane; $a_\lambda$ are the coefficients in the Schur expansion $p=\sum_{\lambda\vdash d}a_\lambda s_\lambda$ of a symmetric homogeneous $p$ of degree $d$; $f^\lambda$ counts standard Young tableaux of shape $\lambda$; and $\Imm_\lambda$ is the immanant of a $d\times d$ matrix attached to the irreducible character $\chi^\lambda$, with $\Imm_{(1^d)}=\det$ and $\Imm_{(d)}=\per$.  Problem 36 asks whether stability forces Schur-positivity outright; Problems 37 and 38 ask for the two endpoint inequalities that would bracket the Schur coefficients between the $\det$ and $\per$ ends of the immanant scale.
\begin{problem}[As posed: Borcea--Br\"and\'en, AIM problem list \citeyearpar{AIMStableProblems}, Problem 36]
Let $p$ be a symmetric, real stable and homogeneous polynomial of degree $d$ in $n$ variables with $d\le n$ and suppose that $p$ has at least one positive (and then all non-negative) coefficients in the monomial basis.  Is $p$ Schur-positive?  That is, when $p(z_1,\ldots,z_n)$ is expanded in the Schur-basis
\[p(z_1,\ldots,z_n)=\sum_{\lambda\vdash d}a_\lambda s_\lambda(z_1,\ldots,z_n),\]
is $a_\lambda\ge0$ for all $\lambda\vdash d$?
\end{problem}
\begin{problem}[As posed: Borcea--Br\"and\'en, AIM problem list \citeyearpar{AIMStableProblems}, Problem 37]
If $A$ is a $d\times d$ matrix then $\det(A)=\Imm_{(1^d)}(A)$ and $\per(A)=\Imm_{(d)}(A)$, where $(1^d)=(1,\ldots,1)$ and $(d)=(1^d)'=(d,0,\ldots,0)$.  Schur's inequality [J,L] asserts that if $A$ is a $d\times d$ positive semi-definite matrix then $\Imm_\lambda(A)\ge f^\lambda\det(A)$ for all $\lambda\vdash d$, where $f^\lambda$ is the number of standard Young tableaux of shape $\lambda$.  A natural real stable extension of this result would be as follows.

\emph{Problem 37.}  Let $p$ be as in Problem 36.  Is $a_\lambda\ge f^\lambda a_{(d)}$ for all $\lambda\vdash d$?
\end{problem}
The answer to both is no already at $d=n=5$.  Let
\begingroup
\scriptsize
\[
J_1=
\begin{pmatrix}
38&19&19&2&36\\
19&38&2&19&36\\
19&2&38&19&36\\
2&19&19&38&36\\
36&36&36&36&72
\end{pmatrix},\quad
J_2=
\begin{pmatrix}
38&19&19&17&21\\
19&38&17&19&21\\
19&17&38&34&36\\
17&19&34&38&36\\
21&21&36&36&42
\end{pmatrix},\quad
J_3=
\begin{pmatrix}
38&34&19&17&36\\
34&38&17&19&36\\
19&17&38&19&21\\
17&19&19&38&21\\
36&36&21&21&42
\end{pmatrix},
\]
\[
J_4=
\begin{pmatrix}
38&4&19&2&21\\
4&8&2&4&6\\
19&2&38&4&21\\
2&4&4&8&6\\
21&6&21&6&42
\end{pmatrix},\qquad
J_5=
\begin{pmatrix}
8&4&4&2&6\\
4&38&2&19&21\\
4&2&8&4&6\\
2&19&4&38&21\\
6&21&6&21&42
\end{pmatrix},
\]
\endgroup
and set
\begin{equation}\label{eq:aim36-q}
q(x_1,\ldots,x_5)=\frac1{648}\det\left(\sum_{i=1}^5x_iJ_i\right).
\end{equation}

\begin{theorem}[\statusfalse: stable Schur positivity, AIM Problem 36]\label{thm:aim36}
The polynomial \eqref{eq:aim36-q} is symmetric, homogeneous of degree five in five variables, real stable, and all $126$ of its ordinary monomial coefficients are strictly positive, yet
\[
[s_{(1^5)}]\,q=-2972<0.
\]
\end{theorem}
\begin{proof}
The leading principal minors of $J_i$ are
\[
\begin{array}{c|rrrrr}
 &\Delta_1&\Delta_2&\Delta_3&\Delta_4&\Delta_5\\ \hline
J_1&38&1083&28728&202176&1119744\\
J_2&38&1083&28728&202176&1119744\\
J_3&38&288&8208&202176&1119744\\
J_4&38&288&8208&46656&1119744\\
J_5&8&288&1728&46656&1119744,
\end{array}
\]
all positive, so Sylvester's criterion gives $J_i\succ0$.  If $z_1,\ldots,z_5$ lie in the open upper half-plane and $(\sum_iz_iJ_i)v=0$ for some $v\ne0$, then
\[
0=\operatorname{Im}\left\langle v,\left(\textstyle\sum_iz_iJ_i\right)v\right\rangle=\sum_i\operatorname{Im}(z_i)\langle v,J_iv\rangle>0,
\]
a contradiction; hence $q$ is real stable, and homogeneity of degree five is immediate.  Exact expansion of the determinant gives
\begin{align*}
q={}&1728m_{(5)}+19440m_{(4,1)}+66555m_{(3,2)}+140910m_{(3,1,1)}\\
&+250440m_{(2,2,1)}+547405m_{(2,1,1,1)}+1182860m_{(1^5)},
\end{align*}
which exhibits both the $S_5$-symmetry and the strict positivity of every monomial coefficient.  Inverting the Kostka matrix over $\mathbb Z$ in the partition order $(5),(4,1),(3,2),(3,1,1),(2,2,1),(2,1,1,1),(1^5)$ yields
\begin{align*}
q={}&1728s_{(5)}+17712s_{(4,1)}+47115s_{(3,2)}+56643s_{(3,1,1)}\\
&+62415s_{(2,2,1)}+56437s_{(2,1,1,1)}-2972s_{(1^5)}.
\end{align*}
The last coefficient alone follows from the single integer identity
\begin{align*}
[s_{(1^5)}]q={}&1728-2(19440)-2(66555)+3(140910)\\
&+3(250440)-4(547405)+1182860=-2972.
\end{align*}
The source package recomputes the Schur expansion a second, independent way, by the bialternant formula, and every step uses rational arithmetic only.
\end{proof}

\begin{remark}[Origin of the pencil]
The matrices are not a numerical accident.  Let $V=\ker B\cong S^{(3,2)}$, where $B$ maps the permutation module on $2$-subsets of $[5]$ to $\mathbb R^5$ by $e_{ab}\mapsto e_a+e_b$, let $P_i$ be the rank-three projection attached to the point stabilizer of $i$, and let $G$ be the Gram matrix of the chosen basis of $V$.  Then $J_i=2G(I+5P_i)$, and the whole one-parameter family $\det(\sum_ix_i(I+tP_i))$ has a closed Schur expansion whose $s_{(1^5)}$ coefficient turns negative for $t$ large; $t=5$ clears denominators.  The standalone write-up \path{counterexamples/aim-problem-36/paper.tex} develops this construction in full.
\end{remark}

\begin{theorem}[\statusfalse: lower endpoint, AIM Problem 37]\label{thm:aim37}
The polynomial \eqref{eq:aim36-q} satisfies the hypotheses of Problem 37 but violates its lower endpoint inequality
\[
a_\lambda\ge f^\lambda a_{(d)}.
\]
\end{theorem}
\begin{proof}
Take $\lambda=(1^5)$, so that $f^{(1^5)}=1$.  The inequality would require $a_{(1^5)}\ge a_{(5)}$, that is, $-2972\ge1728$.
\end{proof}

\begin{remark}[Relation to Problem 38]
The same sign obstruction refutes Problem 38 as well: at $\lambda=(5)$ the upper endpoint would require $1728\le-2972$.  That refutation is not carried here, because Theorem~\ref{thm:pot} refutes Problem 38 by a Schur-\emph{positive} polynomial, which is the stronger statement.
\end{remark}

\section{Macdonald lattice Schur-convexity}
\casecredits{Posed by McSwiggen and Sahi \citeyearpar{McSwiggenSahi2026}; counterexample found by TODO (TODO); formalized by S.~Sra; audited by S.~Sra.}
For parameters $q,t\in(0,1)$ and a partition $\lambda$ with at most $n$ parts, $P_\lambda(x;q,t)$ denotes the Macdonald polynomial in the monic normalization of Macdonald's book, as used in \citet{McSwiggenSahi2026}; at $q=t$ it specializes to the Schur polynomial $s_\lambda(x)$, which is what the refutation below exploits.  The statement concerns two convexity properties of the map $\lambda\mapsto\Omega_\lambda(x;q,t)$ on partitions: \emph{log-convexity}, and \emph{Schur-convexity}, meaning monotonicity with respect to the dominance (majorization) order, so that $\lambda\succeq\mu$ implies $\Omega_\lambda\ge\Omega_\mu$.  The normalization $\Omega_\lambda$ and the lattice $\mathcal L_n$ of evaluation points are defined inside the quoted statement.
\begin{problem}[As posed: C.~McSwiggen and S.~Sahi \citeyearpar{McSwiggenSahi2026}, Theorem 2.1]
For all $q,t\in(0,1)$, $a>0$, and $x\in\mathcal L_n^{q,t,a}$, the function $\lambda\mapsto\Omega_\lambda(x;q,t)$ is log-convex and Schur-convex.

Here, by \S2.1 of the same paper, $\Omega_\lambda(x;q,t)=P_\lambda(x;q,t)/P_\lambda(t^\delta;q,t)$ is the Macdonald polynomial normalized to $1$ at the principal specialization $t^\delta=(t^{n-1},t^{n-2},\ldots,1)$, and the scaled, $t$-shifted, dominant $q$-lattice is
\[\mathcal L_n=\mathcal L_n^{q,t,a}=\left\{(aq^{-\mu_1}t^{n-1},aq^{-\mu_2}t^{n-2},\ldots,aq^{-\mu_n}):(\mu_1\ge\ldots\ge\mu_n)\in\mathbb Z^n\right\}\subset\R^n.\]
\end{problem}
\begin{theorem}[\statusfalse: Theorem 2.1 of \citet{McSwiggenSahi2026}]\label{thm:macdonald}
The stated Schur-convexity fails already in rank two.
\end{theorem}
\begin{proof}
Set $q=t=r\in(0,1)$, $a=1$, and choose lattice index $\nu=(1,0)$.  The resulting lattice point is $x=(1,1)$.  Let $\lambda=(2,0)$ and $\mu=(1,1)$; then $\lambda\succeq\mu$.  At $q=t$, Macdonald polynomials specialize to Schur polynomials, and direct evaluation gives
\[
\Omega_{(2,0)}(1,1;r,r)=\frac{3}{1+r+r^2},\qquad
\Omega_{(1,1)}(1,1;r,r)=\frac1r.
\]
But
\[
1+r+r^2-3r=(1-r)^2>0,
\]
so $\Omega_{(2,0)}<\Omega_{(1,1)}$ for every $0<r<1$, opposite to Schur-convexity.  At $r=1/2$, the values are $12/7<2$.
\end{proof}

\section{Derivatives of the Jacobi-theta kernel}
\casecredits{Posed by Coffey and Csordas \citeyearpar{Csordas2014}; counterexample found by TODO (TODO); formalized by S.~Sra; audited by S.~Sra.}
The kernel $\Phi$ below is the Jacobi theta kernel whose cosine transform is the Riemann $\xi$-function, so inequalities among its derivatives sit directly on the route to the Riemann Hypothesis: the Turán inequalities for the Maclaurin coefficients of $\xi$ are exactly statements of this shape.  A positive function $g$ is \emph{strictly log-concave} when $(g')^2-gg''>0$, so the assertion $J_n(t)>0$ below says precisely that the $(n-1)$st derivative $\Phi^{(n-1)}$ is strictly log-concave, for every order $n$ and every real $t$.
\begin{problem}[As posed: G.~Csordas \citeyearpar{Csordas2014}, Open Problem 4.13, restating Coffey--Csordas Conjecture 2.5]
([5, Conjecture 2.5])  Show that the derivatives of the Jacobi theta function, $\Phi(t)$, are (strictly) log-concave on $\R$; that is, for each $n\in\mathbb N$,
\[J_n(t):=(\Phi^{(n)}(t))^2-\Phi^{(n-1)}(t)\Phi^{(n+1)}(t)>0\quad\text{for }t\in\R.\]
Here, by (4.2) of the same paper, the Jacobi theta function (without the usual factor $4$) is $\Phi(t):=\sum_{n\ge1}\pi n^2\left(2\pi n^2e^{4t}-3\right)\exp\left(5t-\pi n^2e^{4t}\right)$.
\end{problem}
\begin{theorem}[\statusfalse; computer-assisted proof]\label{thm:theta}
At $n=9$ and $t=1/50$, one has $J_9(t)<0$.
\end{theorem}
\begin{proof}
Put $u=\pi m^2e^{4t}$ and define integer polynomials
\[
P_0(u)=2u-3,\qquad P_{k+1}(u)=4uP_k'(u)+(5-4u)P_k(u).
\]
Termwise differentiation gives
\[
\Phi^{(k)}(t)=\sum_{m\ge1}\pi m^2e^{5t-u}P_k(u).
\]
Outward-rounded interval summation, with an explicit geometrically bounded tail, yields
\begin{align*}
\Phi^{(8)}(1/50)&\approx 2.95940368409056694\cdot10^8,\\
\Phi^{(9)}(1/50)&\approx 5.80442287530267756\cdot10^9,\\
\Phi^{(10)}(1/50)&\approx 2.02040519425950258\cdot10^{11},
\end{align*}
and an interval for $J_9(1/50)$ contained strictly in the negative half-line, centered at
\[
-2.61006208371358922\cdot10^{19}.
\]
The source package contains both a high-precision independent evaluator and a Sage \texttt{RealIntervalField} certificate.  The negative margin is many orders of magnitude larger than the certified truncation error.
\end{proof}

\section{Feasible Picard steps for DPP likelihood}
\casecredits{Posed by Mariet and Sra \citeyearpar{MarietSra2015}; counterexample found by TODO (TODO); formalized by S.~Sra; audited by S.~Sra.}
A determinantal point process with positive definite kernel $L$ on a ground set $\mathcal Y$ assigns a subset $Y\subseteq\mathcal Y$ probability proportional to $\det(L_Y)$, the principal minor on $Y$.  Given observed subsets $Y_1,\ldots,Y_n$, learning the kernel means maximizing the log-likelihood
\[
\phi(L)=\sum_{i=1}^n\log\det(L_{Y_i})-n\log\det(I+L),
\]
whose stationarity condition is $\Delta+L^{-1}=L^{-1}$ for the residual
\[
\Delta=\frac1n\sum_{i=1}^nU_i(U_i^*LU_i)^{-1}U_i^*-(I+L)^{-1},
\]
where $U_i$ is the indicator matrix selecting the coordinates in $Y_i$.  This suggests the Picard iteration $L_{k+1}=L_k+aL_k\Delta_kL_k$ of the quoted passage; a step size $a$ is \emph{feasible} when it keeps the iterate positive definite, which the source guarantees for $a\le1/(1-\gamma)$ with $\gamma=\max\{\lambda_{\min}(LZ),1/\lambda_{\max}(I+L)\}\in[0,1]$.  Ascent is proved for $a=1$; the conjecture asks for every feasible $a\ge1$.
\begin{problem}[As posed: Z.~Mariet and S.~Sra \citeyearpar{MarietSra2015}, \S2 (unnumbered)]
Above we showed that for $a=1$ ascent is guaranteed.  Empirically, $a>1$ often works well; Prop.~A.1 presents an easily computable upper bound on feasible $a$.  We conjecture that for all feasible values $a\ge1$, iteration (2.5) is guaranteed to increase the log-likelihood.

The iteration carrying the step size is $L_{k+1}=L_k+aL_k\Delta_kL_k$, displayed as (2.7) in the source; feasibility is the bound $a\le1/(1-\gamma)$ of Prop.~A.1.
\end{problem}
\begin{theorem}[\statusfalse: feasible-step ascent conjecture]\label{thm:dpp}
A feasible step with $a=5$ can strictly decrease the DPP log-likelihood.
\end{theorem}
\begin{proof}
Use ground set $\{1,2\}$ and observations
\[
\mathcal T=\bigl\{\{1\},\{2\},\{1,2\}\bigr\}.
\]
Let
\[
L_0=\begin{pmatrix}3.37200647&2.62325460\\2.62325460&6.07953037\end{pmatrix},
\]
with the displayed decimals interpreted as exact rationals of denominator $10^8$.  For this data,
\[
\Delta(L)=\frac13\left[
\begin{pmatrix}L_{11}^{-1}&0\\0&0\end{pmatrix}
+\begin{pmatrix}0&0\\0&L_{22}^{-1}\end{pmatrix}+L^{-1}\right]-(I+L)^{-1}.
\]
At $a=5$, exact rational arithmetic gives
\[
L_1=L_0+5L_0\Delta(L_0)L_0
\approx\begin{pmatrix}
3.1679125391632645&3.1681701017272776\\
3.1681701017272776&3.4400758121589554
\end{pmatrix}.
\]
Both matrices are positive definite; numerically,
\[
\det L_0\approx13.6187510458,\qquad
\det L_1\approx0.8605575075.
\]
The normalized objective is
\[
\phi(L)=\frac13\bigl(\log L_{11}+\log L_{22}+\log\det L\bigr)-\log\det(I+L).
\]
The inequality $\phi(L_1)<\phi(L_0)$ is equivalent, after exponentiating and cubing, to the exact rational comparison
\[
\frac{(L_1)_{11}(L_1)_{22}\det L_1}{(L_0)_{11}(L_0)_{22}\det L_0}
<\left(\frac{\det(I+L_1)}{\det(I+L_0)}\right)^3.
\]
Numerically, the left and right sides are respectively
\[
0.03359118941299074
\quad\text{and}\quad
0.04354945027745948.
\]
Their exact rational difference is positive.  Hence the step is feasible but descending.
\end{proof}

\section{Variance-sensitive Matrix Spencer}
\casecredits{Posed by A.~S.~Bandeira \citeyearpar{Bandeira2016}; counterexample found by GPT-5.5 (Pro) (TODO); formalized by E.~Akba\c{s} and S.~Sra \citeyearpar{AkbasSra2026}; audited by S.~Sra; contributed by S.~Sra.}
For symmetric matrices $A_1,\ldots,A_n$ the \emph{discrepancy} is $\disc(A_1,\ldots,A_n)=\min_{x\in\{\pm1\}^n}\|\sum_{i=1}^nx_iA_i\|_{\rm op}$.  Spencer's theorem gives an $O(\sqrt n)$ bound in the scalar case, and the Matrix Spencer conjecture asks for the same $O(\sqrt n)$ bound for $n$ symmetric contractions of size $n\times n$.  The variance-sensitive strengthening below replaces the ambient scale $\sqrt n$ by the intrinsic variance parameter $\|\sum_{i=1}^nA_i^2\|_{\rm op}^{1/2}$, which never exceeds $\sqrt n$ and is often far smaller; it is the parameter that controls the norm of a random signing in matrix concentration, and Kyng, Luh and Song established the strengthened bound in the rank-one case.  The dimension regime is the whole point: the conjecture concerns $n$ matrices of size $n\times n$, so a family whose dimension grows faster than the number of summands says nothing about it.
\begin{problem}[As posed: Remark 4.25 of A.~S.~Bandeira's problem collection \citeyearpar{Bandeira2016}, quoted here from its restatement as Conjecture 1.2 of E.~Akba\c{s} and S.~Sra \citeyearpar{AkbasSra2026}]
There exists a universal constant $C$, such that for any choice of $n$ symmetric matrices $A_1,\ldots,A_n\in\R^{n\times n}$ satisfying $\|A_i\|_{\rm op}\le1$ there exists a coloring $x\in\{\pm1\}^n$, such that
\[\left\|\sum_{i=1}^nx_iA_i\right\|_{\rm op}\le C\left\|\sum_{i=1}^nA_i^2\right\|_{\rm op}^{1/2}.\]
\end{problem}
The disproof recorded here is not ours: the family below is Theorem A.1 of \citet{AkbasSra2026}, the first refutation of the conjecture, and this entry reproduces it.

Fix an integer $m\ge2$ and set $n=2^m$, so that the matrices below are $n\times n$ and there are $n$ of them: the dimension is not inflated relative to the number of summands.  Index coordinates by $s=(s_1,\ldots,s_m)\in S:=\{\pm1\}^m$, so $|S|=n$.  Write $p=(1,\ldots,1)$ and $q=(-1,\ldots,-1)$, choose any $U\subseteq S\setminus\{p,q\}$ with $|U|=m$, put $F=S\setminus U$, and fix a bijection $\pi:\{m+1,\ldots,n\}\to F$; note that $p,q\in F$.  Define diagonal matrices
\begin{equation}\label{eq:disc-family}
(A_i)_{s,s}=\frac{s_i}{\sqrt m}\quad(1\le i\le m),
\qquad
A_k=\frac1{\sqrt m}E_{\pi(k),\pi(k)}\quad(m<k\le n),
\end{equation}
where $E_{r,r}$ is the diagonal matrix with a single $1$ in coordinate $r$.  Each $\|A_i\|_{\rm op}=1/\sqrt m\le1$, so the family is admissible.

\begin{theorem}[\statusfalse: variance-sensitive Matrix Spencer]\label{thm:discrepancy}
For the family \eqref{eq:disc-family},
\[
\disc(A_1,\ldots,A_n)=\frac{m-1}{\sqrt m},
\qquad
\left\|\sum_{i=1}^nA_i^2\right\|_{\rm op}^{1/2}=\sqrt{1+\tfrac1m},
\]
so the ratio of the two sides equals $(m-1)/\sqrt{m+1}=\Theta(\sqrt{\log n})$.  No universal constant $C$ can therefore satisfy the variance-sensitive bound, already for $n$ matrices of size $n\times n$.
\end{theorem}
\begin{proof}
Since the matrices are diagonal, so is $\sum_iA_i^2$, and its $s$th entry is
\[
\sum_{i=1}^m\frac{s_i^2}m+\sum_{k=m+1}^n\frac1m\mathbf 1\{s=\pi(k)\}=1+\frac1m\mathbf 1\{s\in F\},
\]
whence $\|\sum_iA_i^2\|_{\rm op}=1+\frac1m$.

For the discrepancy, let $x\in\{\pm1\}^n$ be an arbitrary signing and $X=\sum_{i=1}^nx_iA_i$, so that
\[
X_{s,s}=\frac1{\sqrt m}\sum_{i=1}^mx_is_i+\frac1{\sqrt m}\mathbf 1\{s\in F\}x_{\pi^{-1}(s)}.
\]
Evaluating at the coordinate $s=(x_1,\ldots,x_m)\in S$ makes the first sum equal to $m$, so
\[
\|X\|_{\rm op}\ge|X_{s,s}|\ge\frac{m}{\sqrt m}-\frac1{\sqrt m}=\frac{m-1}{\sqrt m}.
\]
This bound is attained.  Take $x_1=\cdots=x_m=1$ and $x_{k_p}=-1$, $x_{k_q}=1$, where $k_p=\pi^{-1}(p)$ and $k_q=\pi^{-1}(q)$, with the remaining signs arbitrary.  Then $\sqrt mX_{s,s}=\sum_{j=1}^ms_j+\mathbf 1\{s\in F\}x_{\pi^{-1}(s)}$, which equals $m-1$ at $s=p$ and $-(m-1)$ at $s=q$.  For every other $s$ the vector is neither all-plus nor all-minus, so $|\sum_js_j|\le m-2$ and hence $|\sqrt mX_{s,s}|\le m-1$.  Thus $\|X\|_{\rm op}=(m-1)/\sqrt m$, and the infimum over signings is exactly that.

Dividing the two displayed quantities gives $(m-1)/\sqrt{m+1}$, which grows without bound as $m=\log_2n\to\infty$.
\end{proof}

\begin{remark}[Why the dimension regime matters]
An earlier version of this entry used $n$ matrices of size $2^n\times2^n$.  That is the trivial regime: once the dimension is exponential in the number of summands the failure says nothing about Matrix Spencer, whose whole content lies in $d$ comparable to $n$.  The family above is the informative one, with $d=n$.

The family is commutative, and the growth rate $\sqrt{\log n}$ is exactly what the sharp noncommutative Khintchine inequalities predict in the commutative case, so the logarithmic factor is genuine rather than an artifact of the construction.  This leaves open the possibility raised in \citet{AkbasSra2026} that the variance-sensitive bound survives for sufficiently noncommutative families.  The original Matrix Spencer conjecture is untouched: bounding $\|\sum_iA_i^2\|_{\rm op}^{1/2}\le\sqrt n$ recovers only the weaker $O(\sqrt n)$ target, which this family satisfies comfortably.
\end{remark}

\section{A rank-two completely-positive mixed-norm inequality}
\casecredits{Posed by S.~Sra; counterexample found by TODO (TODO); formalized by S.~Sra; audited by S.~Sra.}
For an entrywise nonnegative matrix $A$, define
\[
\|A\|_{p,q}=\left(\sum_j\left(\sum_i |a_{ij}|^p\right)^{q/p}\right)^{1/q}.
\]
The formal rank-two problem asked whether, for $s=6/5$ and $q=6$,
\begin{equation}\label{eq:mixed-target}
\|A\|_{s,q}\le\sqrt{\|A\|_{s,s}\|A\|_{q,q}}
\end{equation}
for every completely positive matrix $A$ of rank at most two.

\begin{theorem}[\statusfalse; computer-assisted exact construction]\label{thm:mixed}
Inequality \eqref{eq:mixed-target} fails for an explicit $21\times21$ completely positive matrix of rank two.
\end{theorem}
\begin{proof}
Let $X$ have 21 rows, grouped as follows.  For each group use the indicated multiplicity $m_g$, weight $w_g$, and direction $y_g$:
\[
\begin{array}{c|cccc}
g&1&2&3&4\\ \hline
m_g&1&18&1&1\\
w_g&46&17&42&1\\
y_g&(1,0)&(1,1/50)&(1,3/50)&(0,1).
\end{array}
\]
Every row in group $g$ is $w_g^{5/6}y_g$, and set $A=XX^{\mathsf T}$.  Thus $A\succeq0$, $A$ is entrywise nonnegative, and $\operatorname{rank}A=2$.  Put
\[
R_A=\frac{\|A\|_{6/5,6}}
{\sqrt{\|A\|_{6/5,6/5}\|A\|_{6,6}}}.
\]
Direct interval evaluation gives
\begin{align*}
R_A&\ge 1.00000061733911154365777590600087863798,\\
R_A&\le 1.00000061733911154365777590600087863800.
\end{align*}
The entire interval lies above one.  The source package regenerates $A$ from the four groups and evaluates the ratio independently at 100 digits; a Sage interval script supplies outward rounding.
\end{proof}

\section{Loewner monotonicity of alternating log-determinants}
\casecredits{Posed by S.~Sra; counterexample found by TODO (TODO); formalized by S.~Sra; audited by S.~Sra.}
For $A_1,\ldots,A_n\succeq0$ and $X\succ0$, define
\[
f(X)=\sum_{\emptyset\ne S\subseteq[n]}(-1)^{|S|+1}
\log\det\left(X+\sum_{i\in S}A_i\right).
\]
The standalone question was whether $f$ is increasing in the Loewner order.  Its derivative is
\[
Df(X)[H]=\tr(C(X)H),\qquad
C(X)=\sum_{\emptyset\ne S}(-1)^{|S|+1}
\left(X+\sum_{i\in S}A_i\right)^{-1}.
\]
Thus monotonicity for all $H\succeq0$ would require $C(X)\succeq0$.

\begin{theorem}[\statusfalse: Loewner monotonicity]\label{thm:logdet}
The function above need not be Loewner-increasing, already for $n=4$ and $3\times3$ positive definite matrices.
\end{theorem}
\begin{proof}
Interpret all displayed six-decimal entries as exact rationals of denominator $10^6$, and take
\[
X=\begin{pmatrix}
0.302766&-0.089057&-0.179370\\
-0.089057&0.027282&0.052935\\
-0.179370&0.052935&0.107617
\end{pmatrix},
\]
\[
A_1=\begin{pmatrix}
0.227979&-1.264394&-0.740843\\
-1.264394&7.052829&4.136214\\
-0.740843&4.136214&2.429309
\end{pmatrix},
\]
\[
A_2=\begin{pmatrix}
712.456619&657.821465&-108.957581\\
657.821465&607.377892&-100.602246\\
-108.957581&-100.602246&16.664150
\end{pmatrix},
\]
\[
A_3=\begin{pmatrix}
2.117478&-0.450289&-2.782586\\
-0.450289&0.096801&0.592006\\
-2.782586&0.592006&3.659335
\end{pmatrix},
\]
\[
A_4=\begin{pmatrix}
0.176290&-0.490538&-0.778637\\
-0.490538&1.373736&2.178964\\
-0.778637&2.178964&3.459701
\end{pmatrix}.
\]
Exact leading-principal-minor checks show that all five matrices are positive definite.  The alternating inverse sum has numerical spectrum
\[
\{-0.511275\ldots,\;499.1699\ldots,\;1110.2784\ldots\}.
\]
More decisively, for the integer vector $v=(844,-249,-475)^{\mathsf T}$, exact rational arithmetic gives
\[
v^{\mathsf T}C(X)v<0.
\]
Taking $H=vv^{\mathsf T}\succeq0$ yields $Df(X)[H]<0$, so $f$ decreases along a positive semidefinite direction.
\end{proof}

\section{Lorentzian Jensen--Bregman metric principle}
\casecredits{Posed by S.~Sra; counterexample found by TODO (TODO); formalized by S.~Sra; audited by S.~Sra.}
A homogeneous polynomial with nonnegative coefficients is \emph{Lorentzian} in the sense of Br\"and\'en and Huh \citeyearpar{BrandenHuh} if all its second-order directional derivatives, taken at positive arguments, are quadratic forms of Lorentzian signature $(+,-,\ldots,-)$; in two variables, a form $\sum_{k}a_kx^ky^{d-k}$ is strictly Lorentzian exactly when the normalized coefficients $a_k/\binom dk$ are strictly log-concave, which is the criterion used below.  Lorentzian polynomials are log-concave on the positive orthant, so $F=-\log G$ is convex there and every midpoint Jensen gap $\J_G$ is nonnegative.  Recall that $d$ is a \emph{semimetric} when it is symmetric, nonnegative, vanishes only on equal arguments, and satisfies the triangle inequality; the triangle inequality is the property that fails below.  In the geometric statement, $K+L=\{k+\ell:k\in K,\ \ell\in L\}$ is the Minkowski sum of convex bodies and $\Vol$ is Lebesgue volume.
\begin{problem}[As posed, paraphrased: S.~Sra, standalone problem; formulation as recorded in \S1 of the case's \path{paper.tex}]
Let $\mathcal C\subset\R^m$ be an open convex cone, let $G:\mathcal C\to(0,\infty)$, and put $F=-\log G$ and
\[\J_G(x,y):=\frac{F(x)+F(y)}2-F\!\left(\frac{x+y}{2}\right).\]
If $G$ is homogeneous and Lorentzian, then $d_G:=\sqrt{\J_G}$ is a semimetric.
\end{problem}
\begin{problem}[As posed, paraphrased: S.~Sra, standalone problem (geometric specialization of the same principle); formulation as recorded in \S1 of the case's \path{paper.tex}]
The log-volume midpoint gap is the square of a distance on convex bodies: that is,
\[d_{\Vol}(K,L):=\left[\log\Vol\!\left(\frac{K+L}{2}\right)-\tfrac12\log\Vol(K)-\tfrac12\log\Vol(L)\right]^{1/2}\]
is a semimetric on convex bodies.
\end{problem}
Both statements fall to a single witness: a strictly Lorentzian bivariate cubic, carried to convex bodies by Shephard's realization theorem.  Throughout, $\J_G(x,y)=\log G(\frac{x+y}2)-\frac12\log G(x)-\frac12\log G(y)$ is the midpoint Jensen gap of $F=-\log G$, and $d_G=\sqrt{\J_G}$.

\begin{theorem}[\statusfalse: Lorentzian Jensen--Bregman metricity]\label{thm:lorentzian}
For
\begin{equation}\label{eq:lorentzian-cubic}
G(x,y)=y^3+\tfrac{11}5xy^2+\tfrac32x^2y+\tfrac1{10}x^3,\qquad(x,y)\in\R_{>0}^2,
\end{equation}
the polynomial $G$ is strictly Lorentzian, yet $d_G=\sqrt{\J_G}$ violates the triangle inequality on $\R_{>0}^2$.
\end{theorem}
\begin{proof}
Writing $G=\sum_{k=0}^3a_kx^ky^{3-k}$ gives $(a_0,a_1,a_2,a_3)=(1,\frac{11}5,\frac32,\frac1{10})$.  By the bivariate criterion \cite[Example~2.3]{BrandenHuh}, strict Lorentzianity is strict log-concavity of $a_k/\binom3k$, and
\[
\left(1,\tfrac{11}{15},\tfrac12,\tfrac1{10}\right)
\text{ satisfies }
\left(\tfrac{11}{15}\right)^2=\tfrac{121}{225}>\tfrac12,\qquad
\left(\tfrac12\right)^2=\tfrac14>\tfrac{11}{150}.
\]
So $G$ is strictly Lorentzian; in particular $F=-\log G$ is convex, and every Jensen gap below is nonnegative.  Take the rational points
\begin{equation}\label{eq:lorentzian-points}
p=(35,2),\qquad q=\left(\tfrac{11}2,\tfrac{23}2\right),\qquad r=\left(\tfrac1{500},15\right).
\end{equation}
With $R_{ab}:=G((a+b)/2)^2/(G(a)G(b))$ and $\J_G(a,b)=\frac12\log R_{ab}$, exact evaluation gives
\[
R_{pq}=\frac{609584616081}{344696430080},\quad
R_{qr}=\frac{1254036986210090216149539001}{1235460453386242764000000000},\quad
R_{pr}=\frac{6141969899088096806341860001}{2794813396007162280000000000}.
\]
Outward-rounded rational enclosures for the three logarithms and their square roots---no floating point---then certify
\[
d_G(p,r)-d_G(p,q)-d_G(q,r)>\frac{7156704176613847159228106087505875849130467}{10^{45}}>0,
\]
so the triangle inequality fails.
\end{proof}

\begin{theorem}[\statusfalse: log-volume distance on convex bodies]\label{thm:volume-distance}
There are full-dimensional convex bodies $K,L\subset\R^3$ and Minkowski combinations $A,B,C$ of them with
\[
d_{\Vol}(A,C)>d_{\Vol}(A,B)+d_{\Vol}(B,C).
\]
Hence the proposed log-volume midpoint gap is not the square of a semimetric.
\end{theorem}
\begin{proof}
For convex bodies in $\R^3$ the relative Steiner polynomial is $\Vol(xK+yL)=\sum_{i=0}^3\binom3iW_i(K,L)x^{3-i}y^i$.  Put
\[
(W_0,W_1,W_2,W_3)=\left(\tfrac1{10},\tfrac12,\tfrac{11}{15},1\right),
\]
which satisfies the full set of Shephard inequalities in dimension three:
\[
W_1^2=\tfrac14>\tfrac{11}{150}=W_0W_2,\qquad
W_1W_2=\tfrac{11}{30}>\tfrac1{10}=W_0W_3,\qquad
W_2^2=\tfrac{121}{225}>\tfrac12=W_1W_3.
\]
By Shephard's realization theorem \citeyearpar[Theorem~4]{Shephard} these are the mixed volumes of some pair $K,L\subset\R^3$, whose Steiner polynomial is then exactly the cubic \eqref{eq:lorentzian-cubic}; the positive endpoint coefficients force $\Vol(K),\Vol(L)>0$, so both bodies are full-dimensional.  Setting $A=p_1K+p_2L$, $B=q_1K+q_2L$, $C=r_1K+r_2L$ for the points \eqref{eq:lorentzian-points}, Minkowski linearity gives $(A+B)/2=\frac{p_1+q_1}2K+\frac{p_2+q_2}2L$ and likewise for the other pairs, so
\[
d_{\Vol}(A,B)^2=\J_G(p,q),\qquad
d_{\Vol}(B,C)^2=\J_G(q,r),\qquad
d_{\Vol}(A,C)^2=\J_G(p,r),
\]
and the violation of Theorem~\ref{thm:lorentzian} transfers verbatim.
\end{proof}

\begin{remark}[What the audit leaves standing]
The midpoint integral identity and the homogeneous primitive decomposition of the earlier positive argument are correct.  What fails is a claimed Gram-determinant identity for the derivative of a Hessian Rayleigh quotient, whose two sides have incompatible homogeneity in the line direction; the omitted calculation therefore conceals a false identity rather than a technical gap.  The full audit, including the exact enclosure procedure, is \path{counterexamples/lorentzian-jensen/paper.tex}.
\end{remark}
